给自定义的层或损失写反向，第一步不是查公式表，是把标量目标的微分整理成一个固定形状：

\[
df=\trace(G^{\trans}dX)=\ip{G}{dX}.
\]

只要式子被摆成这样，\(\nabla_Xf=G\) 就可以直接读出来。上一单元确立了这个读法，本单元提供把式子摆成这个形状的三件工具。

第一件是 trace 的循环性质。它允许把乘积在环上整体旋转，从而把 \(dX\) 转到最右端，剩下的一串就是 \(G^{\trans}\)。二次型 \(f(X)=\trace(X^{\trans}AXB)\) 的梯度 \(AXB+A^{\trans}XB^{\trans}\) 就这样读出来，其中两项对应 \(X\) 在表达式里的两次出现；\(A,B\) 对称时它收缩成 \(2AXB\)，而促成收缩的是对称性，不是二次型这个名目。

第二件和第三件都涉及非多项式的出现方式。\(X^{-1}\) 的微分来自恒等式 \(XX^{-1}=I\)，\(\log\det X\) 的微分来自 Jacobi 公式，两条推导各只有三行。真正需要留心的是它们落到代码上的写法：符号里的 \(X^{-1}\) 表示``解一个线性系统''而不是``构造一个逆矩阵''，符号里的 \(\det X\) 在几百维上根本无法被浮点数表示，必须换成 Cholesky 对角线的对数求和。数学的账和数值的账要分开记，混在一起就会写出能跑但不可信的训练循环。

\subsection{怎么读这一章}

三节按依赖顺序排列。先把二次型练熟——循环移位、转置、路径计数这三样在后两节反复出现。逆矩阵那一节的重点不在公式，而在隐式微分与伴随求解这套写法，它是隐式层、Kalman 滤波和可微物理共用的模式。log-determinant 那一节请把公式段与实现段对照着看，Gaussian 似然和 normalizing flow 的日常成本都压在这两段之间。

每次动用公式之前有一份短对账：矩阵是否可逆，是否对称正定，变量是否受对称或稀疏约束，以及数值实现有没有复用前向已经算出的分解。这四问在本单元的三节里各出现一次。

\subsection{本章篇目}

以下篇目按概念依赖排列；若从具体问题进入，也可以先打开最接近当前困惑的一篇，再沿篇内链接回补。

\begin{itemize}
\tightlist
\item \hyperref[sec:10-Matrix-Calculus/03-Trace-Calculus/01]{``Trace 怎样整理二次型微分''}；
\item \hyperref[sec:10-Matrix-Calculus/03-Trace-Calculus/02]{``逆矩阵求导应写成线性求解''}；
\item \hyperref[sec:10-Matrix-Calculus/03-Trace-Calculus/03]{``行列式与 log-determinant 的微分''}。
\end{itemize}

工具备齐之后，下一单元把它们放到真实模型上：最小二乘、logistic 回归、softmax 交叉熵，以及一张计算图。
